\section{Uniform Sampling over the Feasible Set via the Hit-and-Run Algorithm}
\label{app:hit-and-run}

\subsection{The Feasible Set}

As established in the Computational Strategy section, for each month $t \in \{1,\ldots,Q\}$ we need to construct a set of $I$ feasible asset allocations. An allocation $\Omega=(\omega_1,\ldots,\omega_N)^{\top}$ is feasible for month $t$ if it satisfies three conditions simultaneously: the weights are non-negative ($\omega_i \geq 0$ for all $i$), they sum to one, and the estimated CVaR of the portfolio does not exceed the limit $L(t)$ set by the glidepath.

Formally, the feasible set for month $t$ is defined as:
\begin{equation}
F(t)=\left\{
\Omega \in \mathbb{R}^{N}:
\omega_i \geq 0\ \forall i,\quad
\sum_{i=1}^{N}\omega_i=1,\quad
\cvar(\Omega,t)\leq L(t)
\right\},
\end{equation}
where the estimated CVaR is computed from the $S$ scalar portfolio returns under the scenarios for month $t$,
\begin{equation}
\delta_{s,t} = \Omega^{\top}R_{s,t}, \qquad s=1,\ldots,S,
\end{equation}
by sorting them from lowest to highest and averaging the worst 10\%, since we work with a 90\% confidence level.

Geometrically, $F(t)$ is the intersection of the standard $(N-1)$-simplex with the sublevel set $\{\Omega:\cvar(\Omega,t)\leq L(t)\}$. This latter region has no closed analytical form, since the estimated CVaR is a piecewise linear function of $\Omega$, giving $F(t)$ a convex but irregular geometry. Since direct sampling from $F(t)$ is not straightforward, we use the Hit-and-Run algorithm, a Markov Chain Monte Carlo (MCMC) method designed to sample uniformly over convex bounded sets.

\subsection{The Hit-and-Run Algorithm}

Starting from a feasible point, at each iteration the algorithm picks a random direction, computes the chord of the line through the current point that stays inside the feasible set, and jumps to a point drawn uniformly along that chord. Under standard regularity conditions, the stationary distribution of the chain is the uniform distribution over $F(t)$. We describe each component of the algorithm as implemented.

\subsubsection{Initial Feasible Point}

Initializing the chain requires a point $\Omega^{(0)} \in F(t)$. We attempt the following in order: (i) the equal-weight portfolio $\Omega=(1/N,\ldots,1/N)^{\top}$; (ii) portfolios drawn at random over the simplex; and (iii) if neither of the above satisfies the CVaR constraint, we solve an optimization problem that minimizes the CVaR of $\Omega$ over the simplex using the SLSQP method, and verify that the minimum found is feasible. If the problem is feasible, which in practice always holds if $L(t)$ is not excessively restrictive, $\Omega^{(0)}$ is obtained.

\subsubsection{Direction Generation and Projection onto the Unit-Sum Hyperplane}

Given the current point $\Omega^{(k)}$, a random direction $d$ is generated by drawing $z \sim N(0,I_N)$ and normalizing:
\begin{equation}
d=\frac{z}{\lVert z\rVert}.
\end{equation}
This direction is then projected onto the hyperplane $\{v \in \mathbb{R}^{N}:\sum_i v_i=0\}$ via:
\begin{equation}
d \leftarrow d - \frac{1}{N}\sum_{i=1}^{N}d_i \cdot \mathbf{1},
\end{equation}
and renormalized. This projection is essential: it ensures that moving along $d$ does not alter the sum of the weights, thereby preserving the constraint $\sum_i \omega_i=1$ at every subsequent step.

\subsubsection{Computing the Feasible Line Segment}

The next step is to determine the interval $[t_{\min},t_{\max}]$ of scalar values $\tau$ such that $\Omega^{(k)}+\tau d$ remains in $F(t)$. This interval is computed by combining two sources of constraint.

\paragraph{Non-negativity constraints.}
For each asset $i$ with $d_i \neq 0$, the value of $\tau$ that drives $\omega_i^{(k)}+\tau d_i$ to zero is:
\begin{equation}
\tau_i^{*}=-\frac{\omega_i^{(k)}}{d_i}.
\end{equation}
If $d_i>0$, this yields a lower bound on $\tau$; if $d_i<0$, an upper bound. Taking the maximum of the lower bounds and the minimum of the upper bounds yields the interval $[t_{\min}^{\mathrm{lin}},t_{\max}^{\mathrm{lin}}]$ consistent with non-negativity.

\paragraph{CVaR constraint.}
Within $[t_{\min}^{\mathrm{lin}},t_{\max}^{\mathrm{lin}}]$, the CVaR is a convex function of $\tau$. An adaptive search is carried out from each endpoint of the interval inward to identify the values $t_{\min}^{\cvar}$ and $t_{\max}^{\cvar}$ beyond which $\cvar(\Omega^{(k)}+\tau d,t)>L(t)$. The final feasible segment is:
\begin{equation}
\left[t_{\min},t_{\max}\right]
=
\left[
t_{\min}^{\mathrm{lin}} \vee t_{\min}^{\cvar},
t_{\max}^{\mathrm{lin}} \wedge t_{\max}^{\cvar}
\right].
\end{equation}

\subsubsection{Sampling and Update}

We draw $\tau^{*}\sim\mathrm{Uniform}(t_{\min},t_{\max})$ and set:
\begin{equation}
\Omega^{(k+1)}=\Omega^{(k)}+\tau^{*}d.
\end{equation}
A component-wise clip to $[0,1]$ and renormalization are applied to correct for floating-point rounding errors. If the resulting point satisfies all constraints, it is accepted as the new chain state; otherwise, the chain remains at $\Omega^{(k)}$.

\subsubsection{Burn-in and Sample Collection}

The first $b_{\mathrm{burn}}$ iterations of the chain are discarded (burn-in), as these early draws are still correlated with the starting point $\Omega^{(0)}$ and are not representative of the stationary distribution. We use $b_{\mathrm{burn}}=20$ in our implementation. After burn-in, the subsequent $I$ points are stored as the $I$ feasible allocations for month $t$.

An independent chain is run for each month $t$, so allocations from different months share no internal MCMC state.

\subsection{Handling Serial Correlation Across Months}

As noted in the Computational Strategy section, Hit-and-Run generates the $I$ allocations for each month sequentially. Although the marginal distribution of each draw converges to the uniform distribution over $F(t)$, adjacent draws within a given month exhibit serial correlation, as is inherent in any Markov chain.

If one were to use row $i$ of the $I \times Q$ allocation array directly, that is, the $i$th draw produced by Hit-and-Run in each month, the monthly allocations of portfolio $i$ would not be independent across months. To remove this artificial dependence, we apply an independent random permutation within each column of the array before extracting rows. As a result, the allocation of portfolio $i$ in month $t$ corresponds to a draw selected at random from the $I$ feasible allocations for that month, with no relation to its position in the MCMC sequence.
